\documentclass[11pt,a4paper]{article}

\usepackage[margin=1in]{geometry}
\usepackage{titlesec}
\usepackage{enumitem}
\usepackage{hyperref}
\usepackage{parskip}
\usepackage{graphicx}
\usepackage{float}
\usepackage{booktabs}
\usepackage{amsmath}

\hypersetup{colorlinks=true, linkcolor=blue, urlcolor=blue}

\titleformat{\section}{\large\bfseries}{\thesection}{0.5em}{}
\titleformat{\subsection}{\normalsize\bfseries}{\thesubsection}{0.5em}{}
\setlist{leftmargin=1.4em, topsep=2pt, itemsep=2pt}

\title{\vspace{-1.5cm}\textbf{PKCERT AI \& Software Development Internship \\ Final Capstone: Phishing Email Inspection Desk}}
\author{Abdullah Amir}
\date{AI and Software Development \\ 28 August 2026}

\begin{document}

\maketitle
\vspace{-0.8cm}

Full code lives in this task's directory (\texttt{model\_v2/}, \texttt{backend/},
\texttt{frontend-app/}, \texttt{Dockerfile}, \texttt{tests/}). A live phishing/safe email
classifier: two models trained and evaluated on 18{,}650 real emails (TF-IDF + Logistic
Regression vs.\ a fine-tuned DistilBERT), a FastAPI backend, a React/Framer Motion frontend
built through a structured design workflow, packaged into one Docker image, and deployed live
on Render's free tier. The project's actual thesis -- deploy the small, cheap model and keep
the larger Transformer as a measured comparison, not a live dependency -- is a direct,
deliberate application of a resource-ceiling finding from an earlier task in this internship,
not a fresh discovery.

% ----------------------------------------------------------------
\section*{Part A: Scope, Planning \& Integration}

\textbf{A1 -- problem and why this project.} Phishing email is one of the most common
real-world attack vectors, and most people have no fast way to sanity-check a suspicious
email beyond intuition. The brief required demonstrable integration of prior-task skills
(text processing, embeddings/Transformer modeling, full-stack development) and independent
project scoping; a live phishing classifier satisfies all three while giving the project a
genuine cybersecurity framing.

\textbf{A2 -- what this replaced.} An earlier iteration of this same capstone slot built an
identical architecture around movie-review sentiment classification. It was fully working,
tested, and deployed before being deliberately replaced with this phishing-detection version
-- the underlying engineering pattern (two-model comparison, deploy-the-small-one decision,
Docker/Render pipeline) carried over almost unchanged; the domain, data, models, and entire
frontend were rebuilt.

\textbf{A3 -- tech stack.}

\begin{table}[H]
\centering
\small
\begin{tabular}{@{}p{2.6cm}p{4.3cm}p{7.2cm}@{}}
\toprule
\textbf{Layer} & \textbf{Choice} & \textbf{Why} \\
\midrule
Data & \texttt{zefang-liu/\allowbreak phishing-email-dataset} (HF \texttt{datasets}) & 18{,}650 real emails, permissively hosted, large enough for a meaningful two-model comparison \\[0.3em]
Model A & TF-IDF + Logistic Regression (\texttt{scikit-learn}) & Static-embedding-style features, sub-millisecond inference, deployable within a strict free-tier memory budget \\[0.3em]
Model B & Fine-tuned \texttt{distilbert-base-\allowbreak uncased} (\texttt{transformers}) & Contextual embeddings -- the direct comparison point for Model A; evaluated, never deployed \\[0.3em]
Backend & FastAPI, async, Pydantic validation & Proven pattern reused from an earlier task rather than reinvented \\[0.3em]
Frontend & React + Vite + Framer Motion & Portfolio-grade presentation for this capstone's evaluative audience; builds to static assets, adding no runtime cost to the deploy target \\[0.3em]
Deployment & Docker (Node build stage + Python runtime), Render free tier & Single container; the 512MB memory budget is the binding constraint shaping the whole architecture \\
\bottomrule
\end{tabular}
\end{table}

\textbf{A4 -- hard scope boundary.} Detection only. No UI element, endpoint, or copy
generates, templates, or suggests phishing content, even illustratively -- treated as a
non-negotiable constraint, not a style choice, since this is a defensive-security tool.

% ----------------------------------------------------------------
\section*{Part B: Data Preparation, Model Building \& Evaluation}

\textbf{B1 -- data cleaning, two real issues found.} 19 rows had null/empty text. A second,
less obvious issue surfaced only after reading actual misclassified examples from a first
training run: \textbf{533 rows ($\sim$2.9\%) carried the literal placeholder string
``empty'' as their entire email text} -- a scraping artifact in the source dataset, not real
content, and invisible to a naive null check. Both were dropped. Class balance: 11{,}322
safe vs.\ 7{,}328 phishing ($\sim$61/39) -- handled with \texttt{class\_weight=balanced} and
macro-averaged metrics throughout. Text length is extreme (median 159 words, one
3.5-million-word outlier); raw text is truncated to 20{,}000 characters before any
processing. The dataset ships only a single \texttt{train} split, so an 80/10/10
train/val/test split was created directly, stratified by label.

\textbf{B2 -- Model A: hyperparameter search (validation set).}

\begin{table}[H]
\centering
\small
\begin{tabular}{rrrrr}
\toprule
\textbf{max\_features} & \textbf{ngram\_max} & \textbf{C} & \textbf{Val acc.} & \textbf{Val F1} \\
\midrule
10{,}000 & 1 & 1.0 & 0.9751 & 0.9738 \\
20{,}000 & 1 & 1.0 & 0.9751 & 0.9738 \\
20{,}000 & 2 & 1.0 & 0.9790 & 0.9779 \\
20{,}000 & 2 & 0.5 & 0.9751 & 0.9739 \\
\textbf{20{,}000} & \textbf{2} & \textbf{2.0} & \textbf{0.9807} & \textbf{0.9796} \\
30{,}000 & 2 & 1.0 & 0.9801 & 0.9790 \\
\bottomrule
\end{tabular}
\end{table}

\textbf{B3 -- Model B: a real resource constraint that changed the training plan.} An
initial attempt trained on the full $\sim$14{,}500-row training split at 128 tokens/example,
matching the design used elsewhere in this internship. It was killed after \textbf{12.8
wall-clock hours still short of finishing epoch 1 of 3} -- verified directly via
\texttt{uptime}/\texttt{free -h} as genuine, sustained system-wide contention (load average
consistently above 12 on a 12-core machine, 8GB+ swapped), not a hang or a bug. Rather than
keep waiting on an uncontrollable external constraint, the run was restarted at a leaner
configuration -- a 6{,}000-example stratified subsample, a 96-token cap, 2 epochs -- which
completed in under 40 minutes. This is the same category of decision as the project's
Model A/B deployment choice: adapt the plan to a measured real-world constraint.

\begin{table}[H]
\centering
\small
\begin{tabular}{rrrrr}
\toprule
\textbf{Epoch} & \textbf{Train loss} & \textbf{Train acc.} & \textbf{Val loss} & \textbf{Val acc.} \\
\midrule
1/2 & 0.1695 & 0.9390 & 0.0701 & 0.9768 \\
2/2 & 0.0394 & 0.9890 & 0.0766 & \textbf{0.9785} \\
\bottomrule
\end{tabular}
\end{table}

Validation loss ticks up slightly between epochs while accuracy still edges upward -- the
same overfitting-begins-early signal seen elsewhere in this internship's training runs,
arriving faster here specifically because the training set is smaller. The saved checkpoint
is epoch 2's (best validation accuracy), not epoch 1's.

\textbf{B4 -- the comparison, reported honestly.}

\begin{table}[H]
\centering
\small
\begin{tabular}{lrr}
\toprule
\textbf{Metric} & \textbf{Model A (TF-IDF + LogReg)} & \textbf{Model B (DistilBERT)} \\
\midrule
Test accuracy & \textbf{98.45\%} & 97.73\% \\
Test macro F1 & \textbf{0.9837} & 0.9761 \\
Avg.\ latency & \textbf{0.20ms/example} & 46.09ms/example \\
Artifact size & \textbf{908KB} & 267.9MB \\
Misclassified & 28/1{,}810 (1.5\%) & 41/1{,}810 (2.3\%) \\
Deployed live & \textbf{Yes} & No \\
\bottomrule
\end{tabular}
\end{table}

\textbf{Model A wins on every axis this time, not just size and speed.} Unlike a pattern
this same comparison has produced on other tasks (a Transformer trading size/speed for a
real accuracy gain), Model A's 98.45\% edges out Model B's 97.73\% here -- a fair, disclosed
consequence of the resource-constrained retraining in B3, not evidence that Transformers are
worse at phishing detection in general. Reviewing Model B's own misclassified examples shows
the same error class as Model A's: well-crafted promotional/phishing text that reads as
formally legitimate, and legitimate emails (a masonic-lodge newsletter, a corporate
congratulations note) using unusually formal phrasing. Contextual embeddings did not
resolve this ambiguity class under this training budget -- a narrower, more honest finding
than assuming they always would. This does not change the deployment decision: Model A's
908KB/0.20ms footprint against Model B's 267.9MB/46ms was never a close call.

% ----------------------------------------------------------------
\section*{Part C: Backend \& Frontend Development, Integration}

\textbf{C1 -- API.} FastAPI, async routes, Pydantic validation, structured JSON logging,
CORS. Three endpoints: \texttt{GET /healthz}, \texttt{POST /api/v1/predict}
(\texttt{\{text\}} $\to$ \texttt{\{label, confidence, latency\_ms\}}), \texttt{GET
/api/v1/models} (both models' real measured metrics, read from saved JSON -- Model B's
267.9MB weights are never loaded into this process, only its already-computed results file).
Input validation (1--5{,}000 chars, non-blank) is enforced by Pydantic at the boundary.

\textbf{C2 -- frontend, built through a structured design workflow.} React + Vite + Framer
Motion, built to static assets and served by FastAPI from the same container -- no separate
frontend hosting, no Node process at runtime. Designed via the project's Impeccable
workflow: a recorded product context (\texttt{PRODUCT.md}), a rolled and weighed visual
direction (a \emph{document-authentication-checkpoint} world -- an inspection counter under
raking light, a stamped verdict -- chosen over catalog challengers including a particle-
collider event display and a demoscene terminal aesthetic on audience-identification and
product-clarity axes), a built implementation, a structured finish review, and a recorded
design system (\texttt{DESIGN.md}). The predict tool is the first viewport, not a form under
marketing copy: pasting text triggers a raking-light scan animation and a rubber-stamped
CLEARED/FLAGGED verdict, with confidence rendered as real instrument-panel visual weight
(bar length, brightness, and glow scaled directly off the measured number, never an inert
percentage label).

\textbf{A genuine design-review finding, fixed.} The finish review flagged three uses of a
banned UI pattern (a kicker/eyebrow label above a heading) and an identical scroll-reveal
animation duplicated across four sections. Both were fixed in a follow-up commit -- the
kicker framing folded into heading copy, and each section given a materially different
motion mechanism (a drawn rule line, a scale-snap, no motion at all, a table row-stagger) so
the hero's scan-to-stamp interaction remains the system's one authored motion signature.
Verified resolved on both desktop and mobile screenshots before shipping.

\textbf{C3 -- scope adherence.} Every UI element traces to Part A's in-scope feature list --
no accounts, no history, no multi-class labels, no non-English support, and no interaction
that could read as generating phishing content.

% ----------------------------------------------------------------
\section*{Part D: Testing, Deployment \& Documentation}

\textbf{D1 -- testing.} \texttt{tests/test\_api.py}, 16 cases via \texttt{pytest}: health
check, predict happy-path (2), validation errors (5), model-comparison, frontend/static
serving (2, including resolving the Vite build's content-hashed asset filenames from the
served HTML rather than hardcoding them), and 4 parametrized clear-case label checks. All 16
pass. \textbf{A real bug found before testing even started}: the 533-row ``empty''
placeholder issue (B1), caught by reading misclassified examples, not a data-quality scan.
\textbf{A second real bug, found while building}: the built frontend's Vite output emits
root-relative asset paths (\texttt{/assets/index-xxxx.js}), which silently 404'd against a
fixed \texttt{/static} mount prefix -- fixed by mounting the built directory at \texttt{/}
with \texttt{html=True} instead.

\textbf{D2 -- deployment, verified directly.} Multi-stage Docker build (Node stage builds
the frontend, Python stage runs the backend), non-root user, \texttt{HEALTHCHECK}, final
image \textbf{649MB}. Run with a hard \texttt{--memory=512m} cap -- the exact ceiling that
OOM'd an earlier project in this internship -- the container used \textbf{110MiB (21.48\%)}
of that budget, with all endpoints and the frontend confirmed working through the capped
container itself, not just the uncapped process.

\begin{sloppypar}
\textbf{Live cloud deployment succeeded}, on Render's free tier --
\url{https://day29-phishing-inspector.onrender.com} -- deployed via Render's API with
GitHub \texttt{main}-branch auto-deploy as the CI/CD pipeline. Verified live, not just
built:
\end{sloppypar}

\begin{verbatim}
$ curl https://day29-phishing-inspector.onrender.com/healthz
{"status":"ok","model_loaded":true,"uptime_s":247.55}

$ curl -X POST .../api/v1/predict -d '{"text": "Verify your account ..."}'
{"label":"phishing","confidence":0.994,"latency_ms":9.08}
\end{verbatim}

\texttt{/api/v1/models} and the frontend were confirmed live too, both models' real numbers
rendering correctly in the deployed comparison table.

\textbf{D3 -- documentation.} This report, \texttt{README.md} (setup/usage), Part A's
scoping document, \texttt{PRODUCT.md}/\texttt{DESIGN.md} (the design workflow's product and
visual-system records), and the generated architecture diagram together let a third party
set up and run this project independently.

% ----------------------------------------------------------------
\section*{Part E: Presentation \& Wrap-up}

A 9-slide deck (\texttt{presentation/slides.html}) covers the problem, the Model A/B
decision and why, measured results, the live demo, the deployment-proof certificate,
architecture, engineering challenges, and lessons learned -- styled consistently with the
deployed product's own design system. The live URL above serves as the working demo. Full
internship wrap-up documents (\texttt{DAILY\_LOGS.md}, \texttt{INTERNSHIP\_REPORT.md}) and
the GitHub repository accompany this submission.

\end{document}
